\documentclass{beamer}
\usepackage[T1]{fontenc}
\usepackage[utf8]{inputenc}
\usepackage{textcomp}
\usepackage{eurosym}
\DeclareUnicodeCharacter{00B0}{\textdegree}
\DeclareUnicodeCharacter{20AC}{\euro}
\DeclareUnicodeCharacter{2013}{--}
\DeclareUnicodeCharacter{2014}{---}
\DeclareUnicodeCharacter{2018}{`}
\DeclareUnicodeCharacter{2019}{'}
\DeclareUnicodeCharacter{201C}{``}
\DeclareUnicodeCharacter{201D}{''}
\DeclareUnicodeCharacter{2026}{\ldots{}}
\usepackage[spanish,es-noshorthands]{babel}

% Definición del color corporativo aproximado para Pantone Warm Red C
\definecolor{UMWarmRed}{HTML}{F9423A}

% Tema y esquema de colores
\usetheme{Madrid}
\usecolortheme{rose} % Usamos un esquema de color que pueda combinarse bien

% Configuración de colores usando el color corporativo
\setbeamercolor{title}{bg=UMWarmRed, fg=white}
\setbeamercolor{frametitle}{bg=UMWarmRed, fg=white}
\setbeamercolor{structure}{fg=UMWarmRed}
\setbeamercolor{section number projected}{bg=UMWarmRed,fg=white}
\setbeamercolor{itemize item}{fg=UMWarmRed}
\setbeamercolor{itemize subitem}{fg=UMWarmRed}
\setbeamercolor{enumerate item}{fg=UMWarmRed}

\title[DISEÑO DE EXPERIMENTOS]{METODOLOGÍA Y DISEÑO DE EXPERIMENTOS: \\ CONCEPTOS BÁSICOS}
\author[Alfonso Rosa García]{Alfonso Rosa García }
\institute[UM]{\textbf{Fundamentos del Análisis Económico} }

\date{9 y 11 de marzo de 2026 \\ 10:00--13:00}

\begin{document}

\begin{frame}
  \titlepage
\end{frame}

\begin{frame}{Bibliografía}
    \textbf{Material Fundamental para el Curso}
    \begin{itemize}
        \item Montgomery, D.C. (2019): "Design and Analysis of Experiments", 10th Edition. Wiley.
    \end{itemize}
    
    \textbf{Metodología del Diseño Experimental}
    \begin{itemize}
        \item Kirk, R. E. (2009). "Experimental design", Sage handbook of quantitative methods in psychology, 23-45.        
        \item Tanco-Raínusso, M. (2008). "Metodología para la aplicación del diseño de experimentos (DOE) en la industria." San Sebastián: Universidad de Navarra.
        \item Ruiz de Maya, S. y López-López, I. (2013): “Metodología del diseño experimental”, en Francisco J. Sarabia Sánchez (ed.) Métodos de Investigación Social y de la Empresa (483-499). Madrid: Pirámide.
        \item Ato, Manuel, Juan J. López-García, and Ana Benavente. "Un sistema de clasificación de los diseños de investigación en psicología." Anales de Psicología/Annals of Psychology 29.3 (2013): 1038-1059.        
    \end{itemize}
\end{frame}

\begin{frame}{Referencias}
   \begin{itemize}
        \item Begley, C., Ellis, L. Raise standards for preclinical cancer research. Nature 483, 531–533 (2012). https://doi.org/10.1038/483531a
        \item Camerer CF, Dreber A, Forsell E, Ho TH, Huber J, Johannesson M, et al. (March 2016). "Evaluating replicability of laboratory experiments in economics". Science. 351 (6280): 1433–1436.     
        \item Evans, A. S. (1978). "Causation and disease: a chronological journey: the Thomas Parran lecture". American Journal of Epidemiology, 108(4), 249-258.
        \item Granger, C. W. J., \& Newbold, P. (2014). "Forecasting economic time series". Academic press.
        \item Hill, Bradford. (1965). "The environment and disease: association or causation?". Proceedings of the Royal Society of Medicine, 58, 295-300.    
    \end{itemize}
    
\end{frame}

\begin{frame}{Referencias}

   \begin{itemize}
        \item Ioannidis, John P. A. (2005). "Why Most Published Research Findings Are False". PLOS Medicine. 2 (8): e124.
        \item Open Science Collaboration. (2015). Estimating the reproducibility of psychological science. Science, 349(6251), aac4716.
        \item Prinz, F., Schlange, T., \& Asadullah, K. (2011). Believe it or not: how much can we rely on published data on potential drug targets?. Nature reviews Drug discovery, 10(9), 712-712.
        \item Tröhler, U. (2005). "Lind and scurvy: 1747 to 1795. Journal of the Royal Society of Medicine, 98(11), 519-522.  
   \end{itemize}
   
\end{frame}
   
\end{document}







